\chapter{Recombination, Decoupling, and Acoustic Scales}

Recombination and acoustic scales are introduced at the level needed to connect the background and perturbation sectors to observable angular scales.  Atomic rates and tight-coupling formulas are treated as standard model inputs.

\allowdisplaybreaks

\section{Conditional atomic benchmark}
The benchmark inputs are
\begin{align}
H_0&=70\ \mathrm{km\,s^{-1}\,Mpc^{-1}},\\
\Omega_b&=0.05,
&\Omega_m&=0.30,\\
\Omega_r&=9\times10^{-5},
&T_0&=2.7255\ \mathrm K,\\
Y_p&=0.24.
\end{align}
These values define a reproducible controlled model; they are not an observational fit.

Let
\begin{align}
T(z)&=T_0(1+z),\\
n_H(z)&=n_{H0}(1+z)^3,\\
S(z)&=
\frac1{n_H(z)}
\left(
\frac{m_ek_BT(z)}{2\pi\hbar^2}
\right)^{3/2}
\exp\left[-\frac{\chi_H}{k_BT(z)}\right].
\end{align}
The Saha equilibrium equation is
\begin{equation}
\frac{x_e^2}{1-x_e}=S.
\end{equation}
Multiplication by \(1-x_e\) gives the quadratic equation
\begin{equation}
x_e^2+Sx_e-S=0.
\end{equation}
Hence
\begin{equation}
x_e=\frac{-S\pm\sqrt{S^2+4S}}{2}.
\end{equation}
For \(S>0\), the minus sign is negative, so the physical root is the plus branch.  Rationalizing that branch gives
\begin{align}
x_e(S)
&=\frac{-S+\sqrt{S^2+4S}}2\\
&=\frac{2S}{S+\sqrt{S^2+4S}}\\
&=\frac{2}{1+\sqrt{1+4/S}}.
\end{align}
Thus
\begin{equation}
\boxed{
x_e(S)=\frac{-S+\sqrt{S^2+4S}}2
=\frac{2}{1+\sqrt{1+4/S}}.
}
\end{equation}

\section{Optical depth and visibility}
Ignoring reionization,
\begin{align}
n_e(z)&=x_e(z)n_H(z),\\
\frac{\dd\tau}{\dd z}
&=
\frac{c\sigma_Tn_e(z)}{(1+z)H(z)},\\
\tau(z)&=\int_0^z
\frac{c\sigma_Tn_e(z')}{(1+z')H(z')}
\dd z'.
\end{align}
The survival probability from redshift \(z\) to the present is \(\ee^{-\tau(z)}\).  Its loss per unit redshift is
\begin{align}
-\frac{\dd}{\dd z}\ee^{-\tau(z)}
&=-\left[-\ee^{-\tau(z)}\frac{\dd\tau}{\dd z}\right]\\
&=\ee^{-\tau(z)}\frac{\dd\tau}{\dd z}.
\end{align}
The redshift-space visibility is therefore
\begin{equation}
\boxed{
g_z(z)=\ee^{-\tau(z)}\frac{\dd\tau}{\dd z}.
}
\end{equation}
Integrating the total derivative gives
\begin{align}
\int_0^{z_{\max}}g_z(z)\dd z
&=-\int_0^{z_{\max}}
\frac{\dd}{\dd z}\ee^{-\tau(z)}\dd z\\
&=\ee^{-\tau(0)}-\ee^{-\tau(z_{\max})}\\
&=1-\ee^{-\tau(z_{\max})}.
\end{align}
The controlled Saha benchmark gives
\begin{align}
z_*&=1282.15,\\
z_{\tau=1}&=1289.423097,\\
\Delta z_{\mathrm{FWHM}}&=139.2.
\end{align}
The visibility-normalization residual is
\begin{equation}
\varepsilon_g=4.807365\times10^{-8}.
\end{equation}
Nonequilibrium multilevel recombination and reionization are \text{\normalfont not derived here}.

\section{Baryon loading and sound horizon}
Since nonrelativistic baryons and photons scale as
\begin{equation}
\rho_b(z)=\rho_{b0}(1+z)^3,
\qquad
\rho_\gamma(z)=\rho_{\gamma0}(1+z)^4,
\end{equation}
the loading factor is
\begin{align}
R(z)
&:=\frac{3\rho_b}{4\rho_\gamma}\\
&=\frac{3\rho_{b0}(1+z)^3}
{4\rho_{\gamma0}(1+z)^4}\\
&=\frac{3\Omega_b}{4\Omega_\gamma(1+z)},\\
c_s(z)&:=\frac{c}{\sqrt{3(1+R(z))}}.
\end{align}
The comoving distance travelled by an acoustic signal before decoupling is
\begin{equation}
r_s(t_*)=\int_0^{t_*}\frac{c_s(t)}{a(t)}\dd t.
\end{equation}
Using
\begin{equation}
\frac{\dd z}{\dd t}=-(1+z)H(z),
\qquad
a=(1+z)^{-1},
\end{equation}
one has
\begin{equation}
\frac{\dd t}{a}=-\frac{\dd z}{H(z)}.
\end{equation}
Reversing the limits therefore gives
\begin{equation}
\boxed{
r_s(z_*)
=\int_{z_*}^{\infty}\frac{c_s(z)}{H(z)}\dd z.
}
\end{equation}
The line-of-sight comoving distance is
\begin{equation}
\chi_*:=\int_0^{z_*}\frac{c}{H(z)}\dd z.
\end{equation}
The geometric acoustic multipole is
\begin{equation}
\boxed{
\ell_A:=\pi\frac{\chi_*}{r_s(z_*)}.
}
\end{equation}
For the benchmark,
\begin{align}
r_s(z_*)&=125.784218\ \mathrm{Mpc},\\
\chi_*&=13644.312614\ \mathrm{Mpc},\\
\ell_A&=340.781006.
\end{align}
This \(\ell_A\) is a geometric acoustic ruler, not a claim that the first CMB peak equals \(\ell_A\).

\section{Tight-coupling oscillator}
For constant loading and absent gravitational driving,
\begin{equation}
X''+k^2c_s^2X=0.
\end{equation}
Writing \(\omega_s:=kc_s\), the characteristic roots are \(r=\pm\ii\omega_s\), so
\begin{equation}
X(\eta)=A\cos(\omega_s\Delta\eta)
+B\sin(\omega_s\Delta\eta).
\end{equation}
At the initial time, \(A=X_i\) and \(X_i'=\omega_sB\).  Therefore
\begin{equation}
X(\eta)=X_i\cos(kc_s\Delta\eta)
+\frac{X_i'}{kc_s}\sin(kc_s\Delta\eta).
\end{equation}
For varying \(R(a)\), the acoustic model integrates
\begin{equation}
X''+k^2c_s(a)^2X=0
\end{equation}
through the benchmark background. The midpoint refinement order is
\begin{equation}
p_{\mathrm{acoustic}}=2.021414.
\end{equation}
The final-state error at the finest step is
\begin{equation}
\varepsilon_{\mathrm{acoustic}}
=3.978519\times10^{-5}.
\end{equation}

\section{Diffusion approximation}
Let the positive conformal opacity be
\begin{equation}
\kappa':=a n_e\sigma_T.
\end{equation}
The controlled tight-coupling diffusion estimate is
\begin{equation}
\boxed{
\frac1{k_D^2(\eta_*)}
=
\int_0^{\eta_*}
\frac{\dd\eta}{6\kappa'}
\frac{R^2+\frac{16}{15}(1+R)}{(1+R)^2}.
}
\end{equation}
At leading diffusion order, each Fourier amplitude obeys a local attenuation equation of the form
\begin{equation}
\frac{\dd X_k}{\dd\eta}
=-k^2D_\gamma(\eta)X_k,
\end{equation}
where the tight-coupling coefficient integrates to
\begin{equation}
\int_0^{\eta_*}D_\gamma(\eta)\dd\eta
=\frac1{k_D^2(\eta_*)}.
\end{equation}
Separating variables gives
\begin{align}
\frac{\dd X_k}{X_k}&=-k^2D_\gamma\dd\eta,\\
\log\frac{X_k(\eta_*)}{X_k(0)}
&=-\frac{k^2}{k_D^2},
\end{align}
so the associated damping factor is
\begin{equation}
\mathcal D(k)=\exp\left[-\frac{k^2}{k_D^2}\right].
\end{equation}
The benchmark gives
\begin{align}
r_D&:=k_D^{-1}=4.554719\ \mathrm{Mpc},\\
k_D&=0.219553\ \mathrm{Mpc}^{-1}.
\end{align}
The diffusion formula is a tight-coupling approximation and does not replace a full collision hierarchy.

\section{Open extensions}
\begin{nextsteps}
\item Replace Saha equilibrium with a multilevel nonequilibrium hydrogen-recombination calculation.
\item Add helium recombination and a controlled reionization history.
\item Include neutrino decoupling and its effect on the acoustic background.
\item Derive the full photon collision hierarchy beyond the tight-coupling approximation.
\item Improve diffusion damping and calibrate the drag epoch and sound horizon at precision accuracy.
\end{nextsteps}
